\scanpage{41}
\mspage{128}

Mais
\[
H^0(X,F)\simeq H^0(Y,\underbrace{f_*(F)}_{\substack{\text{loc. const.}\\\text{par hyp}\\\text{c)}}})
\underset{\substack{\searrow\ \text{car }p\\\text{Hot-équiv.}}}{\xrightarrow{\sim}}
H^0(\deleted{\unsure{X}}Y',p^*\underbrace{f_*(F)}_{\deleted{\unsure{f'_*q^*F}}})\sim\ ,
\]
mais
\[
p^*f_*(F)\xrightarrow{\sim} f'_*(q^*(F)),
\]
car $f_*(F)$ et $f'_*(q^*F)$ se calculent ``fibre par
fibre'' par hypothèse. Donc
\[
H^0(X,F)\simeq H^0(Y',f'_*(q^*(F)))\simeq H^0(X',q^*(F)),\quad\text{\unsure{qed.|cqfd.}}
\]

$2^\circ$) Si $F$ est un faisceau en groupes \inserted{loc.} constant sur $X$,
il faut prouver
\[
H^1(X,F)\xrightarrow{\sim}H^1(X',q^{\unsure{*|\wedge}}F)
\]
Le fait que c'est injectif résulte déjà de $1^\circ$,
qui implique que le foncteur $P\longrightarrow \unsure{f|q}^*P$
\[
\underset{\text{sur }X}{F\text{-torseurs}}\longrightarrow
\underset{\text{sur }X'}{q^*(F)\text{-torseurs}}
\]
est pleinement fidèle, par l'argument bien
connu, \unsure{utilisant} les faisceaux $\mathrm{Hom}$. Il
faut voir que ce foncteur est ess. surjectif.

\deleted{La \deleted{fait que} pleine fidélité implique qu'il
suffit de le voir \uline{localement} sur $Y$, \unsure{d'où} donc
\unsure{OPS|OK si} $Y=Y_{/y_0}$ est muni d'un objet final $y_0$.
Mais \deleted{\illegible{1}}\inserted{\unsure{déjà|alors}} $H^1(X,F)\simeq H^1(X_{/y_0},F)\simeq(R^1 f_*F)_{y_0}\simeq H^0(Y,R^1f_*F)$.
D'ailleurs, par hyp. c) $R^1 f_*(F)$ est loc. constant \deleted{et se
calcule fibre par fibre}, donc comme $p$ une
Hot-équivalence, on a
\[
H^0(Y,R^1 f_*F)\simeq H^0(Y',p^*R^1 f_*F),
\]
}

% archive sheet number: 41
